\DocumentMetadata{
  pdfversion=2.0,
  pdfstandard=ua-2,
  lang=en-US,
  tagging=on,
  tagging-setup={math/setup=mathml-SE,tabsorder=structure}
}
\documentclass[11pt,letterpaper]{article}
\usepackage[margin=0.68in,includefoot]{geometry}
\usepackage{newcomputermodern}
\usepackage{amsmath,amscd,mathtools}
\usepackage{tikz,adjustbox}
\providecommand{\square}{\mdlgwhtsquare}
\usepackage[hidelinks]{hyperref}
\hypersetup{
  pdftitle={Numerical Analysis First-Year Exam, May 2025},
  pdfauthor={University of Florida Department of Mathematics},
  pdfsubject={Graduate mathematics examination},
  pdfdisplaydoctitle=true,
  bookmarksopen=true
}
\setcounter{secnumdepth}{0}
\setlength{\parindent}{0pt}
\setlength{\parskip}{0.28em}
\setlength{\emergencystretch}{3em}
\setlength{\leftmargini}{2.15em}
\setlength{\leftmarginii}{2.65em}
\setlength{\leftmarginiii}{3em}
\renewcommand{\labelenumi}{\textbf{\theenumi.}}
\pagestyle{empty}
\raggedbottom
\allowdisplaybreaks
\newenvironment{examproblems}[1]{%
  \begin{enumerate}%
  \setcounter{enumi}{#1}%
  \setlength{\topsep}{0.35em}%
  \setlength{\partopsep}{0pt}%
  \setlength{\itemsep}{0.48em}%
  \setlength{\parsep}{0.18em}%
}{\end{enumerate}}
\newenvironment{examsubparts}{%
  \begin{enumerate}%
  \setlength{\topsep}{0.18em}%
  \setlength{\partopsep}{0pt}%
  \setlength{\itemsep}{0.16em}%
  \setlength{\parsep}{0.1em}%
}{\end{enumerate}}
\newenvironment{examinstructions}{%
  \par\begingroup\itshape
}{%
  \par\endgroup\vspace{0.18em}
}
\makeatletter
\renewcommand\section{\@startsection{section}{1}{0pt}%
  {-1.25ex plus -.25ex minus -.1ex}%
  {0.55ex plus .1ex}%
  {\normalfont\large\bfseries}}
\renewcommand{\@maketitle}{%
  \newpage\null\vskip -1em%
  {\footnotesize\bfseries
    \href{https://www.ufl.edu/}{University of Florida}, \href{https://math.ufl.edu/}{Department of Mathematics}\par}%
  \vskip -0.5em%
  \tagstructbegin{tag=Title}%
    {\LARGE\bfseries\href{https://gma.math.ufl.edu/exams/numerical-analysis-first-year/}{Numerical Analysis First-Year Exam}, May 2025\par}%
  \tagstructend%
  \vskip 0.25em%
  \tagmcbegin{artifact}%
    \hrule height 1pt%
  \tagmcend%
  \vskip 1em%
}
\makeatother
\title{Numerical Analysis First-Year Exam, May 2025}
\author{}
\date{}
\begin{document}
\maketitle
\thispagestyle{empty}
\begin{examinstructions}
Do 4 (four) problems.
\end{examinstructions}
\begin{examproblems}{0}
\item
Consider the function \(g(x)=e^{-x}\).
\begin{examsubparts}
\item[\textnormal{(a)}]
Prove that~\(g\) is a contraction on \(G=[\ln 1.1,\ln 3]\).
\item[\textnormal{(b)}]
Prove that~\(g\) maps \(G=[\ln 1.1,\ln 3]\) into \(G=[\ln 1.1,\ln 3]\).
\item[\textnormal{(c)}]
Prove that \(x_{k+1}=g(x_k)\) converges to an unique fixed point \(z\in G=[\ln 1.1,\ln 3]\) for any initial value \(x_0\in G=[\ln 1.1,\ln 3]\).
\end{examsubparts}
\item
Based on \(u_1(x)=1,u_2(x)=x,u_3(x)=x^2\), use Gram–Schmidt orthogonalization process to compute the three polynomials \(w_1(x),w_2(x),w_3(x)\) which are orthonormal on the interval \([0,1]\) with respect to the inner product \((f,g)=\int_0^1 f(x)g(x)\,dx\). Using these polynomials, find the best approximation in \(P^2[0,1]\) for \(f(x)=x^{\frac12}\).
\item
Consider the finite difference formula 
\[
f'(t_j)=\frac{1}{12h}\left[f(t_j-2h)-8f(t_j-h)+8f(t_j+h)-f(t_j+2h)\right]+O(h^4).
\]
\begin{examsubparts}
\item[\textnormal{(a)}]
Derive this formula by using Taylor’s theorem.
\item[\textnormal{(b)}]
Derive this formula by using Lagrange polynomial representation.
\end{examsubparts}
\item
Assume the numerical quadrature for \(\hat f(\hat x)\) on \([0,1]\) is 
\[
\hat J(\hat f)=\int_0^1 \hat f(\hat x)\,d\hat x\approx \hat Q(\hat f)=\sum_{j=0}^m \hat\alpha_j\hat f(\hat x_j).
\]
 Derive the numerical quadrature of \(J(f)=\int_a^b f(x)\,dx\).
\item
Consider numerically solving the initial value problem 
\[
y'(t)=f(t,y),\quad 0<t\leq t_f,\qquad \text{with }y(0)=\eta.
\]
 Assume~\(f\) is sufficiently differentiable and let~\(h\) denote the step size. Show that all convergent members of the family of methods 
\[
y_{n+2}+(\theta-2)y_{n+1}+(1-\theta)y_n=\frac14h\left[(6+\theta)f_{n+2}+3(\theta-2)f_n\right],
\]
 parameterized by~\(\theta\), are also \(A_0\)-stable.
\end{examproblems}
\end{document}
